\section{Holomorphic Forms and Genus}
\label{sec:holomorphic-forms-and-genus}

Let \(\HX\) be the vector space of global holomorphic one-forms. Define the
analytic genus by \(g(X)=\dim_{\C}\HX\), as in
Definition~\ref{def:analytic-genus}.

\subsection{Finite-dimensionality (analytic route)}

The project's preferred route to \(\dim_{\C}\HX<\infty\) is analytic: a
sup-norm makes \(\HX\) a normed space, Montel's theorem makes the closed
unit ball compact, and Riesz's lemma forces finite dimensionality.

\begin{definition}[Fiber norm on the cotangent bundle]
\label{def:cotangent-fiber-norm}
\lean{JacobianChallenge.Blueprint.cotangentFiberNorm}
For each \(x\in X\) the cotangent fiber
\(T^{*}_{x}X = T_{x}X \to_{\C}\C\) is a one-dimensional complex vector
space. Equip it with the canonical operator norm
\(\|v\|_{x} := \sup_{\xi\in T_{x}X,\ \|\xi\|=1} |v(\xi)|\),
where \(\|\xi\|\) is taken via any chart isomorphism
\(\dd\varphi_{x}: T_{x}X \xrightarrow{\sim} \C\) (any two chart norms
on a one-dimensional fiber differ by a positive scalar, so the operator
norm on \(T^{*}_{x}X\) inherits a chart-independent equivalence class;
to fix a single function rather than an equivalence class we choose
the chart norm coming from the preferred \texttt{chartAt} of the
\texttt{ChartedSpace} structure).

The resulting function
\((x, v) \mapsto \|v\|_{x}\) is continuous on the total space of
\(T^{*}X\): on each chart neighbourhood \(U\) the trivialisation
identifies \(T^{*}_{x}X\) with \(\C\to_{\C}\C\) by an isomorphism
varying continuously in \(x\), and the operator norm is itself
continuous on \(\C\to_{\C}\C\); composition of continuous maps is
continuous. (See Mathlib
\texttt{Bundle.ContinuousLinearMap.vectorBundle} for the vector-bundle
structure on \(T^{*}X\), and the operator-norm continuity from
\texttt{ContinuousLinearMap.opNorm\_le\_iff}.)
\depends{Mathlib operator norm on \(\C\to_{\C}\C\); project-side
\texttt{CotangentSpace} from
\texttt{Jacobian/HolomorphicForms/CotangentBundle.lean}.}
\notready
\end{definition}

\begin{definition}[Sup norm on \(\HX\)]
\label{def:holomorphic-sup-norm}
\uses{def:cotangent-fiber-norm}
\lean{JacobianChallenge.Blueprint.holomorphicSupNorm}
For \(\omega\in\HX\), define
\[
  \|\omega\| := \sup_{x\in X}\|\omega_x\|_{x},
\]
where \(\|\omega_x\|_{x}\) is the fiber norm of
Definition~\ref{def:cotangent-fiber-norm}. Compactness of \(X\) plus
continuity of \(x \mapsto \|\omega_x\|_{x}\) (the composition of the
continuous section \(\omega\) with the continuous fiber-norm function)
implies the sup is attained and finite, by the extreme value theorem
(Mathlib \texttt{IsCompact.exists\_isMaxOn}). The standard checks
\(\|\omega\|=0\iff \omega=0\), \(\|\lambda\omega\|=|\lambda|\|\omega\|\)
and the triangle inequality follow pointwise from the corresponding
properties of the fiber norm; this makes \((\HX,\|\cdot\|)\) a normed
\(\C\)-vector space. (Cf.\ Forster, \emph{Lectures on Riemann
Surfaces}, \S 14.)
\notready
\end{definition}

\begin{lemma}[Local coefficient bound on a chart]
\label{lem:chart-coefficient-bound}
\uses{def:holomorphic-sup-norm}
\lean{JacobianChallenge.Blueprint.chart_coefficient_bound}
Let \(\varphi:U\to V\subseteq\C\) be a chart in the holomorphic atlas
of \(X\). On \(U\) every \(\omega\in\HX\) admits a unique
representation \(\omega = h_{\varphi}(z)\,dz\) with
\(h_{\varphi} : V \to \C\) holomorphic. Then there is a constant
\(C_{\varphi}\ge 0\), depending only on the chart \(\varphi\) (and the
choice of fiber norm on \(T^{*}X\)), such that
\[
  |h_{\varphi}(z)| \;\le\; C_{\varphi}\,\|\omega\|
  \qquad\text{for all } z \in V.
\]
\notready
\end{lemma}

\begin{proofsketch}
Set \(p := \varphi^{-1}(z)\). The chart trivialisation identifies the
fiber \(T^{*}_{p}X\) with \(\C \to_{\C} \C\) isometrically up to a
chart-dependent positive constant \(C_{\varphi}^{-1}\); under this
identification \(\omega_{p}\) corresponds to the linear map
\(\xi\mapsto h_{\varphi}(z)\,\xi\), whose operator norm is exactly
\(|h_{\varphi}(z)|\). Therefore
\(|h_{\varphi}(z)| \le C_{\varphi}\,\|\omega_{p}\|_{p} \le
C_{\varphi}\,\|\omega\|\). The constant \(C_{\varphi}\) is bounded by
the supremum on the (relatively compact) chart image of the
trivialisation comparison factor, which is finite by continuity.
\end{proofsketch}

\begin{lemma}[Montel compactness]
\label{lem:montel-compactness}
\uses{def:holomorphic-sup-norm,lem:chart-coefficient-bound}
\lean{JacobianChallenge.Blueprint.montel_compactness}
The closed unit ball of \(\HX\) is sequentially compact: every sequence
\(\omega_n\in\HX\) with \(\|\omega_n\|\le 1\) for all \(n\) admits a
strictly increasing \(\varphi:\N\to\N\) and a limit
\(\omega_\infty\in\HX\) with \(\|\omega_\infty\|\le 1\) such that
\(\|\omega_{\varphi(n)} - \omega_\infty\| \to 0\) as \(n\to\infty\)
(equivalently, \(\omega_{\varphi(n)} \to \omega_\infty\) uniformly on
\(X\)).
\notready
\end{lemma}

\begin{proofsketch}
\emph{Step 1: finite chart cover.} Compactness of \(X\) supplies a
finite atlas \(\{(U_i,\varphi_i)\}_{i=1}^{N}\) with each
\(\varphi_i(U_i) \supseteq \overline{B(c_i, R_i)}\); choose
\(0 < r_i < R_i\) with \(\bigcup_i \varphi_i^{-1}(B(c_i, r_i)) = X\).

\emph{Step 2: chartwise sup bound.} Lemma~\ref{lem:chart-coefficient-bound}
gives constants \(C_i\) such that the chart coefficient
\(h_{n,i}\) of \(\omega_n\) on \(U_i\) satisfies
\(|h_{n,i}(z)| \le C_i\) for all \(z \in \overline{B(c_i, R_i)}\) and
all \(n\).

\emph{Step 3: Cauchy estimates yield equicontinuity.} For
\(z \in \overline{B(c_i, r_i)}\) the Cauchy integral formula for the
derivative,
\(h_{n,i}'(z) = \frac{1}{2\pi i}\oint_{|w-c_i|=R_i}
\frac{h_{n,i}(w)}{(w-z)^2}\,dw\),
combined with \(|w-z| \ge R_i - r_i\), gives
\(|h_{n,i}'(z)| \le C_i \,R_i / (R_i - r_i)^2\). The family
\(\{h_{n,i}\}\) is therefore uniformly Lipschitz, hence equicontinuous,
on \(\overline{B(c_i, r_i)}\). (Mathlib:
\texttt{DiffContOnCl.circleIntegral\_sub\_inv\_smul} for the
Cauchy formula;
\texttt{circleIntegral.norm\_integral\_le\_of\_norm\_le\_const} for the
\(L^{\infty}\)–to–integral bound.)

\emph{Step 4: Arzelà–Ascoli on each chart.} The Arzelà–Ascoli theorem
(Mathlib \texttt{ArzelaAscoli.isCompact\_closure\_of\_isClosedEmbedding})
extracts a subsequence converging uniformly on
\(\overline{B(c_i, r_i)}\). Iterating across the \(N\) charts
diagonally produces a single subsequence
\(\omega_{\varphi(n)}\) whose chart restrictions all converge
uniformly.

\emph{Step 5: glue and pass to the limit.} The chartwise uniform
limits agree on overlaps (because they are obtained from a single
diagonal subsequence), so they assemble into a function on \(X\) which
is the uniform limit of \(\omega_{\varphi(n)}\). By Weierstrass's
convergence theorem (uniform limit of holomorphic functions is
holomorphic; cf.\ Forster \S 14 or Griffiths–Harris pp.~7–8) the chart
coefficients of the limit are again holomorphic, so the limit is a
holomorphic 1-form \(\omega_\infty\). The sup-norm bound
\(\|\omega_\infty\|\le 1\) follows by taking the limit in
\(|h_{\varphi(n),i}(z)|\le 1\) (with our normalisation; if not unit
norm, by passing to the limit in the explicit estimate).
\end{proofsketch}

\begin{theorem}[Closed unit ball is compact]
\label{thm:hone-unit-ball-compact}
\uses{lem:montel-compactness}
\lean{JacobianChallenge.Blueprint.hone_unit_ball_compact}
Under any normed-\(\C\)-vector-space realisation of \(\HX\) whose norm
agrees with the sup norm of
Definition~\ref{def:holomorphic-sup-norm}, the closed unit ball
\(\overline{B}(0,1) = \{\omega \in \HX : \|\omega\|\le 1\}\) is
(topologically) compact.
\notready
\end{theorem}

\begin{proofsketch}
A normed space is metrisable, hence first-countable; in a first-countable
space sequential compactness and topological compactness coincide
(Mathlib \texttt{UniformSpace.isCompact\_iff\_isSeqCompact} or
\texttt{IsSeqCompact.isCompact}). Apply
Lemma~\ref{lem:montel-compactness} to obtain sequential compactness of
the closed unit ball, then convert. The closed unit ball is closed in
the metric topology by continuity of the norm, so the conclusion is
exactly \texttt{IsCompact (Metric.closedBall 0 1)}.
\end{proofsketch}

\begin{theorem}[Finite-dimensionality from Riesz]
\label{thm:fd-from-riesz}
\uses{thm:hone-unit-ball-compact}
\lean{JacobianChallenge.Blueprint.fd_from_riesz}
A normed \(\C\)-vector space whose closed ball of some positive radius
is compact is finite-dimensional. This is Riesz's lemma; in Mathlib it
is \texttt{FiniteDimensional.of\_isCompact\_closedBall} (file
\texttt{Mathlib/Analysis/Normed/Module/FiniteDimension.lean}). Applied
to a normed-space realisation of \(\HX\) supplied by
Theorem~\ref{thm:hone-unit-ball-compact}, this gives
\(\dim_{\C}\HX<\infty\).
\depends{Mathlib's
\texttt{FiniteDimensional.of\_isCompact\_closedBall}.}
\leanok
\end{theorem}

\begin{proofsketch}
Direct invocation of
\texttt{FiniteDimensional.of\_isCompact\_closedBall ℂ one\_pos h}
where \(h\) is the conclusion of
Theorem~\ref{thm:hone-unit-ball-compact} for the chosen realisation.
\end{proofsketch}

\begin{theorem}[Classical input: Finite-dimensionality (umbrella)]
\label{input:finite-dimensionality}
\uses{thm:fd-from-riesz}
\lean{JacobianChallenge.Blueprint.input_finite_dimensionality}
\notready
For any normed-\(\C\)-vector-space realisation \(H\) of \(\HX\) whose
norm equals the sup norm, \(H\) is finite-dimensional.

This is the umbrella combining the analytic ingredients
(Definitions~\ref{def:cotangent-fiber-norm}--\ref{def:holomorphic-sup-norm},
Lemmas~\ref{lem:chart-coefficient-bound}--\ref{lem:montel-compactness},
Theorems~\ref{thm:hone-unit-ball-compact}--\ref{thm:fd-from-riesz}).
The alternative Riemann-Roch route (\ref{input:riemann-roch}) is named
separately and not used by the analytic finite-dim proof.
\depends{compactness, local coefficient extraction for one-forms,
Montel + Riesz.}
\end{theorem}

\begin{proofsketch}
Combine Theorem~\ref{thm:hone-unit-ball-compact}, which gives
\texttt{IsCompact (Metric.closedBall 0 1)} on the realisation \(H\),
with Theorem~\ref{thm:fd-from-riesz}.
\end{proofsketch}

\begin{theorem}[Finite-dimensionality of holomorphic one-forms]
\label{thm:fd-holomorphic-one-forms}
\uses{input:finite-dimensionality}
\lean{JacobianChallenge.Blueprint.fd_holomorphic_one_forms}
The vector space \(\HX\) is finite-dimensional:
\(\dim_{\C}\HX<\infty\). Equivalently, the analytic genus
\(g(X)=\dim_{\C}\HX\) is a finite natural number.
\notready
\end{theorem}

\begin{proofsketch}
Transport \(\HX\) into the canonical sup-norm realisation supplied by
the project-side \texttt{HolomorphicOneFormBanachData}
(Jacobian/HolomorphicForms/CompactRiemannSurface.lean), apply
Theorem~\ref{input:finite-dimensionality} to conclude
finite-dimensionality of the realisation, and transport back along
the realisation isomorphism. (\(\C\)-linear bijections preserve
finite-dimensionality.)
\end{proofsketch}

\subsection{Riemann-Roch}

We package Riemann-Roch as a single classical input, with named
prerequisite sheaf/cohomology infrastructure that makes the
decomposition visible in the dep graph.

\begin{definition}[Sheaf cohomology of line bundles on \(X\) (placeholder)]
\label{def:sheaf-cohomology-rs}
For an invertible sheaf \(\mathcal{L}\) on \(X\) (equivalently a divisor
class), there are cohomology groups \(H^{0}(X,\mathcal{L})\) and
\(H^{1}(X,\mathcal{L})\), with \(H^{0}(X,\mathcal{O}_X(D))=L(D)\).
\depends{Sheaves on Riemann surfaces, sheaf cohomology, line bundles —
none packaged in Mathlib for compact RS at v4.28.0.}
\notready
\end{definition}

\begin{theorem}[Serre duality on a compact Riemann surface]
\label{thm:serre-duality-rs}
\uses{def:sheaf-cohomology-rs}
For an invertible sheaf \(\mathcal{L}\) and the canonical sheaf \(K\),
\[
  H^{1}(X,\mathcal{L})\cong H^{0}(X,K\otimes\mathcal{L}^{-1})^{\vee}.
\]
\notready
\end{theorem}

\begin{theorem}[Euler characteristic of an invertible sheaf]
\label{thm:euler-char-line-bundle}
\uses{def:sheaf-cohomology-rs,def:divisor-degree}
For an invertible sheaf \(\mathcal{L}=\mathcal{O}_X(D)\),
\(\chi(\mathcal{L}) := h^{0}(\mathcal{L}) - h^{1}(\mathcal{L})
                     = \deg(D) + 1 - g(X).\)
\notready
\end{theorem}

\begin{theorem}[Classical input: Riemann-Roch (umbrella)]
\label{input:riemann-roch}
\uses{input:divisors,thm:serre-duality-rs,thm:euler-char-line-bundle}
\notready
For every divisor \(D\) on \(X\),
\[
  \ell(D)-\ell(K-D)=\deg(D)+1-g(X),
\]
where \(K\) is a canonical divisor and \(\ell(D)=\dim_{\C} L(D)\).
Combines the Euler-characteristic formula
(Theorem~\ref{thm:euler-char-line-bundle}) with Serre duality
(Theorem~\ref{thm:serre-duality-rs}).
\end{theorem}

\begin{proposition}[Genus zero gives a degree-one meromorphic function]
\label{prop:genus-zero-degree-one-map}
\uses{input:riemann-roch}
If \(g(X)=0\), then for every \(p\in X\) there is a meromorphic function
\(f:X\to\CPone\) with exactly one simple pole at \(p\). In particular, \(f\)
has degree one.
\lean{JacobianChallenge.HolomorphicForms.homeomorphic_sphere_of_analyticGenus_eq_zero}
\leanok
\end{proposition}

\begin{proof}
Apply Riemann-Roch to \(D=p\). Since \(\deg(p)=1\) and \(g(X)=0\),
\[
  \ell(p)-\ell(K-p)=2.
\]
The second term is zero because \(K-p\) has negative degree in genus zero.
Thus \(L(p)\) has dimension \(2\), so it contains a nonconstant meromorphic
function with at most a simple pole at \(p\). Such a function cannot have any
other pole and therefore defines a degree-one map to \(\CPone\).
\end{proof}

\subsection{Degree-one holomorphic maps to \texorpdfstring{\(\CPone\)}{CP1}}

\begin{definition}[Branched-cover degree of a nonconstant holomorphic map]
\label{def:branched-degree}
\uses{def:vanishing-order}
For a nonconstant holomorphic \(f:X\to Y\), the degree \(\deg(f)\) is
the cardinality, counted with multiplicity, of any generic fiber. Each
preimage \(p\in f^{-1}(q)\) contributes its ramification index
\(e_p(f)=\ord_p(f - f(p))\ge 1\); the sum is constant in \(q\) and
equals \(\deg(f)\).
\depends{Local power-series order at \(p\) (\texttt{analyticOrderAt});
constancy in the generic fiber from the open-mapping + isolated-zeros
machinery.}
\notready
\end{definition}

\begin{lemma}[Branch locus is finite]
\label{lem:branch-locus-finite}
\uses{def:branched-degree}
The set \(\{p\in X : e_p(f)\ge 2\}\) is finite.
\notready
\end{lemma}

\begin{proofsketch}
\(e_p(f)\ge 2\) iff the local derivative of \(f\) vanishes. The
derivative is itself meromorphic; its zero set is discrete by the
identity theorem, and finite by compactness.
\end{proofsketch}

\begin{lemma}[Degree-one forces no ramification]
\label{lem:degree-one-no-ramification}
\uses{def:branched-degree,lem:branch-locus-finite}
If \(\deg(f)=1\), then \(e_p(f)=1\) for every \(p\in X\), so \(f\) is
unramified.
\notready
\end{lemma}

\begin{proof}
At a generic fiber, the sum of ramification indices equals \(1\), so the
fiber has exactly one point with index \(1\). At a branch point \(p\),
the sum of indices over the fiber is also \(1\), but each index is
\(\ge 1\), so the fiber is a singleton \(\{p\}\) with \(e_p(f)=1\).
\end{proof}

\begin{theorem}[Local biholomorphism from unramified holomorphic]
\label{thm:local-biholo-unramified}
\uses{lem:degree-one-no-ramification}
An unramified nonconstant holomorphic map \(f:X\to Y\) is a local
biholomorphism: every \(p\in X\) has a neighborhood mapped
biholomorphically onto a neighborhood of \(f(p)\).
\depends{Mathlib's local inverse function theorem on \(\C\) plus the
fact that unramification means the local derivative is nonzero.}
\notready
\end{theorem}

\begin{theorem}[Degree-one is bijective]
\label{thm:degree-one-bijective}
\uses{def:branched-degree,thm:local-biholo-unramified}
A degree-one holomorphic map \(f:X\to Y\) is bijective.
\notready
\end{theorem}

\begin{proofsketch}
Surjective: a degree-\(d\) holomorphic map between compact Riemann
surfaces is surjective (the image is open by open-mapping and closed by
compactness). Injective: by
Lemma~\ref{lem:degree-one-no-ramification}, every fiber has exactly one
point with ramification index 1, hence cardinality 1.
\end{proofsketch}

\begin{theorem}[Compact-Hausdorff bijective continuous map is a homeomorphism]
\label{thm:compact-bijection-homeo}
A continuous bijection between compact Hausdorff spaces is a
homeomorphism.
\depends{Mathlib's \texttt{Continuous.homeomorphOfCompactToT2OfBijective}
or analogous topology fact (PRESENT).}
\leanok
\end{theorem}

\begin{theorem}[Classical input: Degree-one maps are isomorphisms (umbrella)]
\label{input:degree-one-isomorphism}
\uses{thm:degree-one-bijective,thm:local-biholo-unramified,thm:compact-bijection-homeo}
\notready
A nonconstant holomorphic map \(f:X\to\CPone\) of degree one is a
biholomorphism.
\end{theorem}

\begin{proof}[Proof of Theorem~\ref{thm:genus-zero-classification}]
If \(g(X)=0\), Proposition~\ref{prop:genus-zero-degree-one-map} and
Input~\ref{input:degree-one-isomorphism} identify \(X\) biholomorphically,
hence homeomorphically, with \(\CPone\), whose underlying topological space is
\(\Sphere\).

Conversely, if \(X\) is homeomorphic to \(\Sphere\), the topological genus is
zero. The comparison theorem between analytic genus
\(\dim_{\C}H^{0}(X,\Omega^{1}_{X})\) and topological genus gives \(g(X)=0\).
\end{proof}

\begin{formalizationnote}
The current Lean target splits this into the two directions in
\code{Jacobian/HolomorphicForms/GenusZeroClassification.lean}. The forward
direction needs Riemann-Roch, meromorphic functions, and degree-one map
classification. The reverse direction needs comparison of analytic and
topological genus.
\end{formalizationnote}
